% ==============================================================
\section{関連研究}\label{sec:related}
% ==============================================================

\paragraph{イベント帰属パイプライン.}
頻出パターンのサポート変動を外部イベントに帰属させるフレームワーク~\cite{dsaa2025}は、
変化点検出・帰属スコアリング・置換検定・BH補正・重複排除の5ステップで構成される。
本研究はこの1次元パイプラインを多次元に拡張し、
空間的近接度と時空間置換検定を導入する。

\paragraph{空間スキャン統計.}
Kulldorff~\cite{kulldorff1997spatial}の空間スキャン統計は、
点過程における空間クラスターを検出する。
時空間拡張~\cite{kulldorff2001prospective}は円筒形窓を使用する。
これらはアイテムセットパターンを扱わず、
検出された密集が\emph{なぜ}生じたか（イベントへの帰属）には答えない。
本研究はパターンサポートの変動をイベントに帰属させる点で根本的に異なる。

\paragraph{密部分テンソルマイニング.}
DenseAlert~\cite{shin2017densealert}、M-Zoom~\cite{shin2018mzoom}は
多次元テンソルにおける密ブロックを検出する。
これらは密度\emph{異常}を検出するが、
検出された密度変動と外部イベントの統計的関連を検定する機能を持たない。

\paragraph{時空間パターンマイニング.}
軌跡パターン~\cite{giannotti2007trajectory}やコロケーションパターン~\cite{shekhar2001colocation}は
空間的パターンを扱うが、パターンのサポート\emph{変動}の原因帰属は対象としない。

\paragraph{計量経済学のイベント・スタディ法.}
MacKinlay~\cite{mackinlay1997event}のイベント・スタディ法は
企業イベントの株価への影響を評価する。
単一のアウトカム変数を対象としており、
数千のパターン時系列を同時に扱う本研究の設定とは異なる。
また空間次元は考慮されない。

\paragraph{サロゲートデータと置換検定.}
Lancaster et al.~\cite{lancaster2018surrogate}は
時系列データに対するサロゲートデータ生成法を体系的に整理している。
本研究の時空間置換検定は、
時間的自己相関を円形シフトで保存するサロゲート法を
空間構造の保存に拡張したものである。
